\begin{initiativkarte}{RobotLab HAW Hamburg}

  % Bild (optional)
  \IfFileExists{bilder/robotfussball.jpg}{%
    \begin{center}
      \includegraphics[width=\linewidth,height=5cm,keepaspectratio]{bilder/robotfussball.jpg}
    \end{center}
    \vspace{0.4cm}
  }{}

  % Kurzbeschreibung
  \textit{\textcolor{hawgray}{Robotfußball ist ein Forschungs- und Lehrprojekt an der HAW Hamburg, bei dem autonome Roboter Fußball spielen und dabei Methoden aus KI, Robotik und Multi-Agenten-Systemen eingesetzt werden.}}

  \vspace{0.5cm}
  \hawrule

  % Beschreibung
  \textbf{\textcolor{hawblue}{Was ist Robotfußball?}}\\[0.3cm]

  Robotfußball ist eine Forschungs- und Lernplattform, die zeigt, wie Informatik und Robotik in der Praxis zusammenwirken.

  Im Mittelpunkt stehen mobile autonome Roboter, die in einem Spiel gegeneinander antreten. Dabei werden reale Probleme aus der Robotik behandelt, insbesondere:

  \begin{itemize}
      \item Navigation mobiler Roboter
      \item Multi-Agenten-Systeme
      \item Reaktive Systeme
      \item Künstliche Intelligenz
  \end{itemize}

  Robotfußball verbindet Forschung und Praxis und zeigt, dass Informatik nicht nur abstrakt ist, sondern direkt in realen Anwendungen sichtbar wird.

  Internationale Wettbewerbe wie der RoboCup und die FIRA (Federation of International Robot-Soccer Association) bilden den Rahmen für diese Disziplin.

  \vspace{0.5cm}

  \textbf{\textcolor{hawblue}{Robotfußball an der HAW}}\\[0.3cm]

  An der HAW Hamburg wird Robotfußball im Rahmen des RobotLabs und des Wahlpflichtprojekts „Mobile Roboter“ seit vielen Jahren eingesetzt.

  Studierende arbeiten hier mit echten Robotersystemen und entwickeln eigene Strategien für Wahrnehmung, Steuerung und Teamverhalten.

  \vspace{0.3cm}

  \textbf{Verwendete Hardware und Systeme:}

  \begin{itemize}
      \item AKSEN-Board als zentrale Steuerungseinheit
      \item Infrarot-Dioden zur Wahrnehmung
      \item Sharp Distanzsensoren für Abstandsmessung
      \item Taster und Schalter für Eingaben
      \item Motoren und Aktoren zur Bewegung
  \end{itemize}

  Die Roboterplattform wird meist mit LEGO-basierten Chassis aufgebaut und flexibel erweitert.

  \vspace{0.5cm}

  \textbf{\textcolor{hawblue}{Wissenschaftlicher Hintergrund}}\\[0.3cm]

  Neben dem spielerischen Aspekt dient Robotfußball als Grundlage für Forschung in Bereichen wie:

  \begin{itemize}
      \item Regelungstechnik
      \item digitale Bildverarbeitung
      \item künstliche Intelligenz
      \item autonome Systeme
      \item industrielle Automatisierung
  \end{itemize}

  \vspace{0.5cm}

  % Tags
  \tag{Robotik}
  \tag{KI}
  \tag{Forschung}
  \tag{Lehre}
  \tag{Autonome Systeme}

  \vspace{0.6cm}
  \hawrule

  % Infozeilen
  \infozeile{\faUsers}{Zielgruppe: Studierende der Informatik und Ingenieurwissenschaften}
  \infozeile{\faGraduationCap}{Semester: Ab höheren Semestern besonders geeignet}
  \infozeile{\faMapMarker*}{Ort: RobotLab HAW Hamburg}
  \infozeile{\faGlobe}{Website: \href{https://users.informatik.haw-hamburg.de/~kvl/somsem05/peters/einleitung.html}{Robotfußball Einführung}}

\end{initiativkarte}